% !TEX encoding = UTF-8
% !TEX program = lualatex

\documentclass[lualatex,ja=standard]{bxjsarticle}

% --- パッケージ設定 ---
\usepackage{graphicx}
\usepackage{listings}
\usepackage{xcolor}
\usepackage{enumitem}
\usepackage{url}
% hyperrefは最後に読み込み、[hidelinks]で赤枠を消します
\usepackage[hidelinks,pdfpagemode=UseNone]{hyperref}

% --- コードブロックのデザイン設定 ---
\lstset{
    basicstyle=\ttfamily\small,
    frame=single,
    breaklines=true,
    backgroundcolor=\color{gray!10},
    rulecolor=\color{gray!30},
    numbers=left,
    numberstyle=\tiny\color{gray},
    showstringspaces=false
}

\title{Mac用　Windows仮想環境構築マニュアル}
\author{作成者　福山翔太}
\date{\today}

\begin{document}

\maketitle
\tableofcontents 
\newpage

\section{初めに}
これは、macOS上でWindowsアプリケーションを利用するために、Mac内にWindows仮想環境を構築するマニュアルです。
VMware Fusionというソフトを使います。

\section{構築手順}
\subsection{VMware Fusionのインストール}
\begin{enumerate}[label=\arabic*.]
    \item 公式サイト（\url{https://support.broadcom.com}）にアクセスし、左のメニューのlog inからアカウントを作成する。
    \item ログイン後、右上の雲をクリックし、VMware Cloud Foundationを選ぶ。
    \item ページ左上のMENUから My Downloads を選ぶ。
    \item 青い枠内のFree Software Downloads available HERE をクリックする。
    \item 検索欄にVMware Fusionと入力し、同じ名前のものを選択する。
    \item 一番上のバージョンを選ぶ。(作成時は25H2)
    \item 25H2u1を選ぶ。
    \item 規約を開いてチェックを入れる。
    \item 雲のマークからダウンロードする。
    \item ダウンロードした .dmgファイルをダブルクリックし、表示されたアイコンをダブルクリック。
    \item 指示に従ってインストール完了。
\end{enumerate}
＊この時点では仮想マシンはMacに認識されていないため、画面間のカーソル移動がカクついたり、コピー＆ペーストができないので注意。

\subsection{Windowsのインストール}
\begin{enumerate}[label=\arabic*.]
    \item MicrosoftからWindowsを取得を選択する。
    \item ダウンロードが完了したら続けるをクリックし、指示通り進む。
    \item ファームウェアはUEFIを選び、暗号化パスワードを決める。
    \item 新しい仮想ディスクを作成を選ぶ。
    \item 設定の詳細が表示され、終了をクリック。
    \item 保存場所を選択し、保存。
    \item 黒い画面に大きな再生ボタンが出たらクリックし、指示通り進む。
    \item セットアップオプションでは「Windowsのインストール」を選択する。
    \item プロダクトキーがありませんをクリックする。
    \item Windows Home を選択。
    \item Windows起動後、メニューバーの「仮想マシン」から「VMware Tools のインストール」を選択する。
    \item エクスプローラーでDVDドライブ内の \texttt{setup.exe} を実行する。
    \item インストール完了後、Windowsを再起動する。
\end{enumerate}

\section{ChromatoPROダウンロード}
\begin{enumerate}[label=\arabic*.]
    \item 仮想マシン内のEdgeで、以下のURLにアクセスし、.NET Framework 2.0 SP2をダウンロードする。\\
    % URL内の特殊文字エラーを防ぐためverbを使用
    \verb|https://www.microsoft.com/ja-jp/download/details.aspx?id=9834|
    
    \item 検索バーに「cmd」と入力し、コマンドプロンプトを「管理者として実行」する。
    
    \item 黒い画面に以下のコマンドを貼り付ける。\\
    \verb|DISM /Online /Enable-Feature /FeatureName:NetFx3 /All /LimitAccess /Source:D:\sources\sxs|
    
    \item 「正常に終了しました」と表示されたら、ChromatoProのインストールが可能になる。
\end{enumerate}

\end{document}